\documentclass[11pt]{article}
\usepackage[letterpaper,margin=1in]{geometry}
\usepackage[T1]{fontenc}
\usepackage{lmodern,microtype}
\usepackage{amsmath,amssymb,amsthm,mathtools}
\usepackage{booktabs,tabularx,array}
\usepackage{xcolor,hyperref}
\definecolor{inkblue}{RGB}{30,61,92}
\hypersetup{colorlinks=true,linkcolor=inkblue,citecolor=inkblue,
 urlcolor=inkblue,pdftitle={The Theta Transform in Three Dimensions},
 pdfauthor={Drafted for Edward Baker}}
\newtheorem{proposition}{Proposition}[section]
\newtheorem{lemma}[proposition]{Lemma}
\numberwithin{equation}{section}
\newcommand{\R}{\mathbb R}
\newcommand{\C}{\mathbb C}
\newcommand{\dd}{\,\mathrm d}
\newcommand{\T}{\mathcal T}
\newcommand{\Sph}{\mathbb S^2}
\newcommand{\Sch}{\mathcal S}
\setlength{\emergencystretch}{2em}
\setlength{\parskip}{3pt}
\title{The Theta Transform in Three Dimensions\\[4pt]
 \large Density, geometry, and observation}
\author{Drafted for Edward Baker}
\date{28 October 2026}
\begin{document}
\maketitle

\begin{abstract}
The Gaussian coefficients of the theta source identify a weighted
three-dimensional radial Laplacian. Starting from this identity, we
realize the source as a positive spatial density and recover the complete function \(H(t,z)\) by an explicit
integral transform. The density is independent of the spectral variable
\(z\); all spectral dependence belongs to the kernel. Logarithmic radius
identifies the transform as a one-dimensional Fourier transform after
angular integration. This gives a precise account of the information
retained by the observation, the information it forgets, and a restricted
inverse when radiality is prescribed. We then characterize the densities that preserve every spectral value
under this working kernel. Its choice remains open for further
investigation; no auxiliary dynamics is selected.
\end{abstract}

\section{Why the theta source leads to three dimensions}\label{sec:source}

The spatial formulation begins with a differential identity already
present in the theta source \cite[Note01, Sections~2--3]{Archive}.
In the normalization of this investigation,
\begin{equation}\label{eq:originalsource}
 \Phi(u)=\sum_{n\ge1}
 \bigl(2\pi^2n^4e^{9u}-3\pi n^2e^{5u}\bigr)
 e^{-\pi n^2e^{4u}}.
\end{equation}
The function \(H(t,z)\) is the cosine observation of
\(e^{tu^2}\Phi(u)\) on \(u\ge0\). We seek a spatial description
that retains this complete source and explains its differential structure.

\subsection{The Gaussian radius and the Laplacian}

The exponential in \eqref{eq:originalsource} suggests
\(r=e^{2u}\): each factor becomes the radial Gaussian
\(e^{-\pi n^2r^2}\). For a radial profile in dimension \(d\),
\begin{equation}\label{eq:gaussianlaplacian}
 \Delta_d=\partial_r^2+\frac{d-1}{r}\partial_r,
 \qquad
 \Delta_d e^{-\pi n^2r^2}
 =\bigl(4\pi^2n^4r^2-2d\pi n^2\bigr)e^{-\pi n^2r^2}.
\end{equation}
Consequently
\[
 \tfrac12r^{5/2}\Delta_d e^{-\pi n^2r^2}
 =\bigl(2\pi^2n^4r^{9/2}-d\pi n^2r^{5/2}\bigr)e^{-\pi n^2r^2}.
\]
The coefficient \(-3\pi n^2\) in the source selects \(d=3\)
within this radial Gaussian realization. Define
\begin{align}
 A(r)&=\sum_{n\ge1}e^{-\pi n^2r^2},\qquad
 \Theta(r)=1+2A(r),\label{eq:theta}\\
 Q(r)&=\Delta_3\Theta(r)
   =4\pi\sum_{n\ge1}n^2(2\pi n^2r^2-3)e^{-\pi n^2r^2}.
 \label{eq:Q}
\end{align}
The series and its derivatives converge locally uniformly for \(r>0\),
so summing the Gaussian identity gives
\begin{equation}\label{eq:radialmotivation}
 \boxed{\displaystyle
 \Phi(u)=\tfrac12r^{5/2}\Delta_3 A(r)
        =\tfrac14r^{5/2}Q(r),\qquad r=e^{2u}.}
\end{equation}
Thus three dimensions arise from the original differential coefficients.
The summation index remains \(n\in\mathbb N\); introducing a
three-dimensional lattice would change the source.

This is also the earlier primitive identity in spatial coordinates.
For \(\psi(u)=r^{1/2}A(r)\), the relation
\(\partial_u=2r\partial_r\) gives
\begin{equation}\label{eq:primitiveconjugation}
 (\partial_u^2-1)\psi(u)
 =4r^{5/2}\left(A''(r)+\frac2rA'(r)\right)=8\Phi(u).
\end{equation}
The source coordinate \(u\) is a logarithmic radius, not itself a
Euclidean radius. Equation~\eqref{eq:primitiveconjugation} is the
operator identity motivating the spatial definitions below.

\subsection{Reflection becomes spherical inversion}

The Jacobi identity at zero argument is
\begin{equation}\label{eq:jacobi}
 \Theta(r)=r^{-1}\Theta(1/r).
\end{equation}
It follows from the modular transformation of \(\theta_3\) with
lattice parameter \(ir^2\) \cite[Eq.~20.7.32]{DLMF}.

The three-dimensional radial Laplacian makes the reflection especially
simple. Direct differentiation gives
\[
 \left(\partial_r^2+\frac2r\partial_r\right)
        \big[r^{-1}f(1/r)\big]
 =r^{-5}\left(f''(1/r)+2r f'(1/r)\right).
\]
Applying this identity to \eqref{eq:jacobi} yields
\begin{equation}\label{eq:Qreflection}
 Q(r)=r^{-5}Q(1/r).
\end{equation}
Every summand of \(Q\) is positive for \(r\ge1\), since
\(2\pi n^2r^2-3>0\). Reflection therefore gives \(Q(r)>0\) for every \(r>0\).


Here \(Q=\Delta_3\Theta\) is a classical identity on \(r>0\).
There is a point source if the Laplacian is taken on all of \(\R^3\):
modularity gives \(\Theta(r)=r^{-1}\) plus a remainder flat at zero,
and hence, distributionally,
\begin{equation}\label{eq:pointsource}
 \Delta_{\R^3}\Theta=Q(|X|)-4\pi\delta_0.
\end{equation}
The smooth function \(Q\) is the regular part used in the observation;
we do not substitute the full distributional Laplacian for it.

Set
\begin{equation}\label{eq:source}
 q_t(u)=e^{tu^2}\Phi(u),\qquad u\in\R.
\end{equation}
Reflection exactly cancels the power \(r^{5/2}\) in
\eqref{eq:radialmotivation}, so
\begin{equation}\label{eq:even}
 \Phi(-u)=\Phi(u),\qquad q_t(-u)=q_t(u).
\end{equation}
The original half-line observation consequently has the two equivalent
forms
\begin{equation}\label{eq:originalH}
 H(t,z)=\int_0^\infty q_t(u)\cos(zu)\dd u
       =\frac12\int_{\R}q_t(u)e^{izu}\dd u.
\end{equation}
The factor \(1/2\) is essential: the whole line contains both reflected
halves of the original source.

\begin{lemma}[Decay]\label{lem:decay}
For every real \(t\), \(q_t\) is an even, positive Schwartz function.
For every \(B\ge0\) and integer \(k\ge0\),
\[
 \int_{\R}|u|^k e^{B|u|}q_t(u)\dd u<\infty.
\]
The same domination holds for the derivatives needed below, uniformly
when \(t\) ranges over a compact real interval.
\end{lemma}
\begin{proof}
For \(r\ge1\), differentiated theta sums are bounded by a polynomial
in \(r\), times \(e^{-c r^2}\), for some \(c>0\).
Indeed, for each fixed integer \(m\),
\(\sum_{n\ge1}n^m e^{-\pi n^2r^2}\le C_m e^{-\pi r^2}\);
finite differentiation changes only the polynomial factors.
Substituting \(r=e^{2u}\) gives domination of the form
\(C\exp(C' u+T u^2-c e^{4u})\) on \(u\ge0\), including
each fixed derivative. Reflection supplies the other half-line.
Smoothness at \(u=0\) follows from the smooth identity
\eqref{eq:Qreflection}.
\end{proof}

\subsection{What this suggests about the kernel}

Substituting \(u=\tfrac12\log r\) and \(\dd u=\dd r/(2r)\)
in the whole-line observation yields
\begin{equation}\label{eq:radialobservation}
 H(t,z)=\frac1{16}\int_0^\infty
 e^{t\log^2r/4}Q(r)\,r^{3/2}\,r^{iz/2}\dd r.
\end{equation}
The oscillatory factor is therefore the Fourier character of logarithmic
radius, \(r^{iz/2}=e^{iz\log r/2}\). With Euclidean volume
\(\dd X=r^2\dd r\dd\omega\), the remaining radial weight leads to
the positive spatial density derived in Section~\ref{sec:lift}.
This is the motivation for the working kernel
\(K(z,X)=e^{iz\ell(X)}\), where \(\ell(X)=\tfrac12\log|X|\).

The Laplacian identity motivates the geometry without selecting a unique
observation kernel or an auxiliary evolution. The exponential kernel
keeps the original Fourier character visible; the cosine kernel agrees
on the initial even source but retains less information on general
later densities, as Section~\ref{sec:information} explains. The choice
remains open as the kernel and evolution are considered together.
In particular, the occurrence of \(\Delta_3\) in the source identity
is not a heat equation for the density. The physical parameter \(t\)
multiplies the source by \(e^{tu^2}\); an auxiliary parameter \(s\)
would require its own law preserving the chosen observations.

\section{The objects and their meanings}\label{sec:objects}

The basic relation will be
\begin{equation}\label{eq:master}
 \boxed{\displaystyle
 H(t,z)=\int_{\R^3}\rho(t;X)\,K(z,X)\dd X.}
\end{equation}
Here \(X\) is a spatial variable and \(z\) is a spectral variable.
The density \(\rho\) describes where the source is placed.
The kernel \(K\) describes how that source is observed.
The scalar \(H\) is the result of the observation.

For real \(t\), the canonical density constructed below is positive.
Cancellation occurs in the oscillatory kernel. Thus the signed density
from the earlier formulation is replaced by a density--kernel pairing:
the factor carrying \(z\) is displayed, not absorbed into \(\rho\).

\begin{center}
\renewcommand{\arraystretch}{1.15}
\begin{tabularx}{\textwidth}{@{}>{$}l<{$}X@{}}
\toprule
\text{Symbol}&Meaning\\
\midrule
t & Newman deformation parameter; changing it changes the source.\\
z\in\C & Spectral argument of \(H\). On the real axis we write \(z=a\).\\
u\in\R & Logarithmic source coordinate.\\
X=(x,y,w),\ r=|X| & Auxiliary position and its Euclidean radius.\\
\ell(X) & The source label read from a spatial point.\\
\rho(t;X) & Spatial density, independent of \(z\).\\
j(t;X) & Spatial gradient \(\nabla_X\rho(t;X)\), a vector field.\\
s & A possible later parameter indexing densities, distinct from \(t\); the kernel stays fixed.\\
\bottomrule
\end{tabularx}
\end{center}

The third Cartesian coordinate is written \(w\), so it cannot be confused
with the spectral variable \(z\). Spatial differentiation always means
\[
 \nabla_X=(\partial_x,\partial_y,\partial_w),
 \qquad \Delta_X=\partial_x^2+\partial_y^2+\partial_w^2.
\]
A conservation flux would require a separate definition.
Here \(j\) means the gradient only; no additional capital \(J\) is used.

\section{A spatial realization without spectral dependence}\label{sec:lift}

Write \(X=r\omega\), with \(\omega\in\Sph\), and let
\(\dd\sigma=\dd\omega/(4\pi)\) be normalized area measure on the sphere.
Our source label is
\begin{equation}\label{eq:phase}
 \ell(X)=\frac12\log|X|,\qquad X\ne0.
\end{equation}
The map \(X=e^{2u}\omega\) identifies punctured three-dimensional
space with \(\R\times\Sph\). Its volume element is
\begin{equation}\label{eq:jacobian}
 \dd X=r^2\dd r\dd\omega
       =8\pi r^3\dd u\dd\sigma(\omega).
\end{equation}
We distribute the measure \(\tfrac12q_t(u)\dd u\) uniformly over the
angular variables. Equations \eqref{eq:radialmotivation}, \eqref{eq:source}, and
\eqref{eq:jacobian} then force the density
\begin{equation}\label{eq:rho}
 \boxed{\displaystyle
 \rho(t;X)=\frac{q_t(\ell(X))}{16\pi |X|^3}
 =\frac{e^{t\log^2r/4}Q(r)}{64\pi\sqrt r},\qquad r=|X|>0.}
\end{equation}
Equivalently, the entire construction is the measure identity
\begin{equation}\label{eq:measure}
 \rho(t;X)\dd X=\tfrac12q_t(u)\dd u\dd\sigma(\omega).
\end{equation}
Thus the angular variables initially carry no additional source data.
They provide a spatial realization of the original logarithmic coordinate.

\begin{proposition}[The spatial transform]\label{prop:transform}
With
\begin{equation}\label{eq:kernel}
 K(z,X)=e^{iz\ell(X)},
\end{equation}
the density \eqref{eq:rho} satisfies \eqref{eq:master} for every real
\(t\) and complex \(z\). It extends by \(\rho(t;0)=0\) to a smooth
Schwartz function on \(\R^3\), positive away from the origin.
Its total mass is \(H(t,0)\).
\end{proposition}
\begin{proof}
Integrating \eqref{eq:measure} over the sphere gives
\[
 \int_{\R^3}\rho(t;X)e^{iz\ell(X)}\dd X
 =\tfrac12\int_{\R}q_t(u)e^{izu}\dd u=H(t,z).
\]
Absolute convergence follows from Lemma~\ref{lem:decay}.
At large \(r\), the density and each spatial derivative have Gaussian
decay after polynomial and logarithmic factors are absorbed.
Equations \eqref{eq:even} and \eqref{eq:rho} give
\begin{equation}\label{eq:rhokelvin}
 \rho(t;r)=r^{-6}\rho(t;1/r).
\end{equation}
Near zero, the Gaussian at \(1/r\) dominates every inverse power of
\(r\), including all factors produced by differentiation. Every
Cartesian derivative therefore extends to zero with value zero.
\end{proof}

The phase \(\ell\) is singular at the origin, but the density is not.
The value of \(K\) at that one point may be assigned arbitrarily;
it does not change any integral. For complex \(z\), the magnitude
of the kernel is \(r^{-\operatorname{Im}z/2}\), which is dominated
by the density at both ends.

This symmetric realization uses the full space.
The regions \(r>1\) and \(r<1\) each contribute half of the
\emph{cosine} observation. The earlier exterior-only realization used
\(q_t(\ell)/(8\pi r^3)\) on \(r\ge1\) and zero inside.
Here its exterior density is halved before the reflected interior
is included. Thus no copy of \(H\) has been added twice, and no jump
has been introduced at the unit sphere.

The requested gradient notation is now unambiguous:
\begin{equation}\label{eq:gradient}
 j(t;X):=\nabla_X\rho(t;X),\qquad
 j(t;X)=\partial_r\rho(t;r)\frac Xr\quad(r>0),\qquad j(t;0)=0.
\end{equation}
It measures spatial steepness. In particular,
\(\nabla_X\cdot j=\Delta_X\rho\) is a differential identity,
not an evolution equation.

\section{What the transform sees}\label{sec:information}

To understand the forward map, fix \(t\) and temporarily write
\(\rho(X)\) for any integrable spatial density.
Let
\[
 \overline\rho(r)=\int_{\Sph}\rho(r\omega)\dd\sigma(\omega),
 \qquad
 g_\rho(u)=8\pi e^{6u}\overline\rho(e^{2u}).
\]
Then \(g_\rho(u)\dd u\) is the distribution of logarithmic radius,
and
\begin{equation}\label{eq:factorization}
 (\T\rho)(z)
 :=\int_{\R^3}\rho(X)e^{iz\ell(X)}\dd X
 =\int_{\R}g_\rho(u)e^{izu}\dd u
\end{equation}
whenever the integral converges. For the canonical density,
\(g_\rho=q_t/2\).

This factorization describes the information content:
\[
 \text{spatial density}
 \ \longrightarrow\
 \text{logarithmic-radius distribution}
 \ \longrightarrow\
 \text{spectral function}.
\]
The first arrow averages over angles. The second is a Fourier transform
with the sign convention displayed in \eqref{eq:factorization}.
In radial variables the kernel is \(r^{iz/2}\), so the same observation
is also a Mellin-type transform.

\subsection{Angular information is lost}

A function with zero spherical average at every radius is invisible to
\(\T\). This failure of injectivity persists among positive smooth
densities. For example, if \(\rho\) is the canonical density, then
\begin{equation}\label{eq:angularambiguity}
 \rho_\varepsilon(r\omega)
 =\rho(r)\big[1+\varepsilon(\omega_1^2-\omega_2^2)\big],
 \qquad |\varepsilon|<1,
\end{equation}
has the same transform for every \(z\).
The multiplier is positive, and the added angular term has mean zero.
Flatness of \(\rho\) at the origin gives a smooth extension there.
Thus the forward map is not invertible on unrestricted three-dimensional
densities. This is an explicit ambiguity, not an inference merely from
a comparison of dimensions.

\subsection{A restricted inverse}

If \(g_\rho\in\Sch(\R)\), Fourier inversion gives
\begin{equation}\label{eq:inverse}
 g_\rho(u)=\frac1{2\pi}\int_{\R}(\T\rho)(a)e^{-iau}\dd a.
\end{equation}
This is the standard Schwartz-space inversion theorem
\cite[Theorem~5.66]{Hunter}, with the forward normalization changed
to the one in \eqref{eq:factorization}.
Consequently the observed function determines the spherical average
\(\overline\rho\). If radiality is prescribed, it determines the full
density, for \(r>0\), by
\begin{equation}\label{eq:radialinverse}
 \rho(r)=\frac1{16\pi^2r^3}
        \int_{\R}(\T\rho)(a)e^{-ia\log r/2}\dd a.
\end{equation}
For the canonical source, evenness reduces this to
\begin{equation}\label{eq:canonicalinverse}
 \rho(t;r)=\frac1{8\pi^2r^3}
       \int_0^\infty H(t,a)\cos\!\left(\frac a2\log r\right)\dd a.
\end{equation}
There is therefore a natural inverse on a specified radial class.
It is not an inverse for the angular freedom in
\eqref{eq:angularambiguity}, and it does not recover an unknown later
spatial deformation.

\subsection{Why the kernel choice matters}

For the canonical density, reflection \(u\mapsto-u\) cancels the sine
part of \(e^{izu}\). Hence we may equally write
\[
 H(t,z)=\int_{\R^3}\rho(t;X)\cos(z\ell(X))\dd X.
\]
For a general density these two kernel families carry different
information. The exponential kernel recovers all of \(g_\rho\);
the cosine kernel recovers only its even part
\[
 g_\rho^{\rm even}(u)=\tfrac12\big(g_\rho(u)+g_\rho(-u)\big).
\]
Thus retaining the exponential representation fixes more data than
retaining only the cosine observation on an unrestricted class.
The two choices agree on the canonical even source, but their
constraints on a future representation should not be conflated.

\section{Spectral differentiation and collisions}\label{sec:collisions}

All spectral dependence is now explicit in the kernel.
Lemma~\ref{lem:decay} permits differentiation under the integral,
uniformly on compact subsets of \(\C\):
\begin{equation}\label{eq:spectraljets}
 \partial_z^k H(t,z)
   =\int_{\R^3}(i\ell(X))^k\rho(t;X)e^{iz\ell(X)}\dd X.
\end{equation}
In particular \(H(t,\cdot)\) is entire.
The physical parameter has the separate meaning
\(\partial_t\rho=\ell^2\rho\), so the same integral gives
\begin{equation}\label{eq:newman}
 H_t=-H_{zz}.
\end{equation}
No spatial evolution is used to derive this identity.

At a real observation \(z=a\), evenness of \(q_t\) gives the real formulas
\begin{align}
 H(t,a)&=\int_{\R^3}\rho(t;X)\cos(a\ell(X))\dd X,\label{eq:realH}\\
 H_a(t,a)&=-\int_{\R^3}\ell(X)\rho(t;X)
                         \sin(a\ell(X))\dd X.\label{eq:realderivative}
\end{align}
A collision at \((T,A)\) means that both integrals vanish.
Equivalently, the positive measure \(\rho(T;X)\dd X\) has zero pairing
with the two test functions
\[
 \cos(A\ell(X)),\qquad -\ell(X)\sin(A\ell(X)).
\]
These are two global cancellation conditions.
They do not specify a spatial point at which the density vanishes.
Nor is the spatial gradient \(j=\nabla_X\rho\) the spectral derivative
of \(H\): the latter differentiates the kernel, as
\eqref{eq:spectraljets} shows.

A multiple zero in the complex spectral plane has the analogous
conditions \(H=H_z=0\); a zero of higher multiplicity imposes
further vanishing pairings in \eqref{eq:spectraljets}.
This is a transform description of the event, not an exclusion mechanism.

\section{The invariant pairing with a fixed kernel}
\label{sec:representations}

The transform is fixed before any auxiliary evolution is chosen.
Let \(s\) index densities \(\rho_s(t;X)\), with
\(\rho_0(t;X)=\rho(t;X)\), and define
\[
 j_s(t;X)=\nabla_X\rho_s(t;X).
\]
The spatial coordinates, volume measure, and kernel remain the same:
\begin{equation}\label{eq:fixedkernel}
 K(z,X)=e^{iz\ell(X)},\qquad \ell(X)=\tfrac12\log|X|.
\end{equation}
Neither \(\rho_s\) nor \(j_s\) contains a spectral argument, and
\(K\) contains no auxiliary parameter \(s\).

The defining requirement for an admissible family is
\begin{equation}\label{eq:invariant}
 \boxed{\displaystyle
 \int_{\R^3}\rho_s(t;X)K(z,X)\dd X=H(t,z)
 \quad\hbox{for every }z\hbox{ and every admitted }s.}
\end{equation}
The same forward integral then recovers \(H\) from the current density
at any \(s\). An inverse transform is unnecessary for this recovery.
We require the integrability needed to make the observation meaningful;
for an entire spectral function, a convenient condition is
\[
 \int_{\R^3}e^{B|\ell(X)|}|\rho_s(t;X)|\dd X<\infty
 \qquad\hbox{for every }B\ge0.
\]

The quantifier \emph{for every \(z\)} is substantive.
Since \(K(0,X)=1\), conservation of total mass preserves \(H(t,0)\).
It does not by itself preserve the other spectral values.
The full invariant is a family of pairings with a fixed family of
kernels. Fixing the kernel defines the observation map; requiring its
output to stay constant is an additional condition.

There are three distinct levels of preservation:
\[
 \rho_s=\rho
 \quad\Longrightarrow\quad
 \T\rho_s=\T\rho
 \quad\Longrightarrow\quad
 \int_{\R^3}\rho_s\dd X=\int_{\R^3}\rho\dd X.
\]
The first fixes the spatial integrand pointwise, the second fixes every
spectral observation, and the third fixes only total mass.
The middle condition is the one imposed here.

\subsection{Exactly what is fixed}

Fourier uniqueness and \eqref{eq:factorization} make
\eqref{eq:invariant} equivalent to
\begin{equation}\label{eq:marginalinvariant}
 g_{\rho_s(t;\cdot)}(u)=\tfrac12q_t(u)
 \quad\hbox{for almost every }u,
\end{equation}
or, equivalently,
\begin{equation}\label{eq:sphericalinvariant}
 \overline{\rho_s}(t;r)=\rho(t;r)
 \quad\hbox{for almost every }r>0.
\end{equation}
For integrable densities, uniqueness follows, for example, by convolving
the difference of the marginals with a Gaussian and taking an
approximate-identity limit. Equality on the real spectral axis suffices;
the exponential moments justify the complex observations as well.

Thus preserving every integral recovers a pointwise statement
\emph{after angular integration}: the one-dimensional integrand
\(g_{\rho_s}(u)e^{izu}=\tfrac12q_t(u)e^{izu}\) is fixed for almost
every \(u\). The three-dimensional integrand
\(\rho_s(X)K(z,X)\) may still change.
This explains how an integral-preserving requirement can retain the
radial information that a pointwise requirement was meant to protect.
It follows from the information seen by this particular kernel.

If a later density is also radial,
\eqref{eq:sphericalinvariant} forces it to equal the initial density.
This is a consequence of the fixed observation map, not a dynamical
symmetry theorem.

\subsection{The nonradial freedom}

Fixing a spherical average does not fix the density at every angle.
Since the canonical density is positive for \(r>0\), the same constraint
can be expressed as
\begin{equation}\label{eq:angularfreedom}
 \rho_s(t;r\omega)=\rho(t;r)A_s(t;r,\omega),\qquad
 \int_{\Sph}A_s(t;r,\omega)\dd\sigma(\omega)=1.
\end{equation}
For nonnegative densities, \(A_s\ge0\); signed densities can be
admitted with the stated absolute-integrability condition.
Equation~\eqref{eq:angularfreedom} is a description of the allowed
representations, not a chosen evolution equation.
The example \eqref{eq:angularambiguity} exhibits this freedom explicitly.

A nonradial density can change at a given spatial point while the
average over that sphere stays fixed. Changes at different radii may
also be coupled, provided each spherical average is retained.
The observation does not specify how those changes should occur.
For smooth members the gradient records their geometry through
\[
 j_s=A_s\,\nabla_X\rho+\rho\,\nabla_X A_s.
\]
In particular, an angularly varying \(A_s\) produces tangential
components of \(j_s\), even though the initial gradient is radial.

The same information can be stated without coordinates. Let
\(\nu_{s,t}\) be the distribution of the fixed label \(\ell\):
\[
 \nu_{s,t}(B)=\int_{\ell^{-1}(B)}\rho_s(t;X)\dd X,
 \qquad B\subset\R\ \hbox{Borel}.
\]
Then the invariant transform is exactly the assertion
\[
 \dd\nu_{s,t}(u)=\tfrac12q_t(u)\dd u.
\]
For nonnegative densities this fixed marginal also supplies every
required exponential moment. The forward transform is therefore
recoverable at every admitted \(s\), while the unresolved angular
information remains in the spatial density.

This formulation fixes the density, gradient, kernel, and conserved
observation before any auxiliary equation is proposed. Choosing a
useful dynamics inside this class is a separate problem.

\section*{Acknowledgments and scope}

Drafted for Edward Baker with substantial assistance from OpenAI's
GPT-6-astra in Codex, at the recorded reasoning effort \emph{max},
28 October 2026. Edward Baker proposed separating the density from the
spectral kernel, clarifying the transformed space and its possible
inverse, keeping the kernel independent of the auxiliary parameter,
and developing the setup before selecting auxiliary flows.
The Laplacian motivation was restored at his request on 28 October 2026
with assistance from GPT-6-astra in Codex, recorded reasoning effort
\emph{xhigh} (chat0014).
The aspiration to a Book exposition is an aspiration to clarity and
economy, in the sense associated with Paul Erd\H{o}s.

The theta normalization and reflection are inherited from the investigation;
Fourier inversion is classical. The calculations here organize those
ingredients into a transform formulation. Internal checking does not
constitute independent specialist review. No new zero-exclusion or
de Bruijn--Newman bound is asserted.

\begin{thebibliography}{9}
\bibitem{DLMF}
NIST Digital Library of Mathematical Functions,
Chapter~20, especially
\href{https://dlmf.nist.gov/20.2.E3}{Eq.~20.2.3} and
\href{https://dlmf.nist.gov/20.7.E32}{Eq.~20.7.32}.
The zero-argument transformation is used in \eqref{eq:jacobi}.

\bibitem{Hunter}
J.~K. Hunter, \emph{Notes on Partial Differential Equations},
revision 18 June 2014, Section~5.B.2, Theorem~5.66.
\href{https://www.math.ucdavis.edu/~hunter/pdes/pde_notes.pdf}
{Author's UC Davis notes}.
Our forward Fourier integral has no prefactor; its inverse has
the factor \(1/(2\pi)\) in one dimension.

\bibitem{Archive}
E. Baker research archive, prepared with LLM assistance,
\emph{Conservative Theta Flows and Collision Amplification},
28 October 2026, Section~2, and Investigation05 Notes01 and25.
This manuscript retains that definition of \(H\), while using the
symmetric full-space density \eqref{eq:rho}.
\end{thebibliography}
\end{document}
